\section{Security Requirements Catalogue}
\label{sec:catalogue}

This section refines the stakeholder requirements R1--R9 into a catalogue
of verifiable system security requirements. The measurable
non-functional requirements of the first version of the plan are now
part of this catalogue.

\subsection{Method}

The requirements follow the risk process of NIST SP~800-160~\cite{nist800160v1}.
The team identified the assets (device key and \ac{PIN}, tokens, account
data and money, transfer orders, audit log, \ac{HSM} keys and release
artefacts), drew the level-1 \ac{DFD} of Annex~A and applied \ac{STRIDE}
to each of its elements~\cite{shostack2014}. This gave the 20 threats of
Annex~B. Each threat is answered by at least one requirement and each
requirement traces back to at least one threat, so there are no orphans.

Each requirement has one ``shall'', names the actor and uses numbers
instead of adjectives~\cite{hooks1993}. Table~\ref{tab:req} gives its
type, functional~(F), non-functional~(N) or assurance~(As), and whether
it applies to the system~(Sys), the Customer~(Cus) or the Admin~(Adm).
The \emph{Trace} column links it to the threats of Annex~B and to R1--R9,
and the verification uses the tools of Table~\ref{tab:lifecycle}.

\subsection{Coverage and Traceability}

The catalogue has 23 requirements: five for confidentiality, five for
integrity, four for availability, eight for authentication,
authorisation and accountability, and one assurance requirement. Four
apply to users, \mbox{REQ-05}, \mbox{REQ-19} and \mbox{REQ-22} to the Customer and \mbox{REQ-18} to the
Admin, and the other 19 apply to the system. Every threat of Annex~B is
answered, and T01 (stolen \ac{PIN}) is answered at five layers. \ac{PSD2} maps to \mbox{REQ-06} and \mbox{REQ-15}, \ac{GDPR} to \mbox{REQ-02} to
\mbox{REQ-04}, \ac{DORA} to \mbox{REQ-14} and \mbox{REQ-21}, and the \ac{CRA} essential
requirements to \mbox{REQ-01}, \mbox{REQ-02}, \mbox{REQ-09} and \mbox{REQ-10}.

\subsection{Conflicts and Residual Risk}

Short tokens (\mbox{REQ-16}) work against psychological acceptability, so tokens
are renewed silently while the user is active. A strict lockout would let
an attacker lock customers out, so \mbox{REQ-13} blocks only one device and
keeps recovery open, and \mbox{REQ-12} limits devices instead of IP addresses.
The main residual risk is a deceived customer who authorises a transfer
to a mule account with every factor present. \mbox{REQ-05}, \mbox{REQ-20} and \mbox{REQ-22} reduce it
but cannot remove it, so it is recorded in the residual-risk list signed
by the PO.
